% Course 1 research-paper starter.
% This is a generic teaching template, NOT a CVPR or NeurIPS submission template.
% Switch to the official venue template only after checking current instructions.
\documentclass[11pt]{article}

\usepackage[margin=1in]{geometry}
\usepackage{graphicx}
\usepackage{amsmath,amssymb}
\usepackage{booktabs}
\usepackage{hyperref}
\usepackage[numbers]{natbib}

\title{A Precise, Evidence-Based Research Question}
\author{Student Author(s)\\
\small Course 1: Deep Learning and Neural Computing Foundations}
\date{\today}

\begin{document}
\maketitle

\begin{abstract}
State the question, method, principal evidence, and bounded conclusion.
Write this section after the results are stable. Do not include invented
or placeholder results in a submitted manuscript.
\end{abstract}

\section{Introduction}
What problem are we studying, why does it matter, and what precise question
will this paper answer?

\section{Related Work}
Summarize relevant work critically. Explain what is known and what remains
uncertain. Cite original sources and verify every reference.

\section{Research Question and Hypothesis}
State the hypothesis in a form that could be contradicted by evidence.

\section{Method}
Describe the model, data provenance, preprocessing, split, baseline,
intervention, controls, metrics, seeds, and compute budget.

\section{Experimental Setup}
Give enough detail for another researcher to reproduce the experiment.
Explain any deviations from the planned protocol.

\section{Results}
Report all relevant runs, appropriate uncertainty, and clearly labeled
figures and tables. Do not report only the best run.

\begin{table}[ht]
  \centering
  \caption{Replace this placeholder with measured results before submission.}
  \label{tab:results}
  \begin{tabular}{lcc}
    \toprule
    Method & Metric (higher/lower is better) & Compute cost \\
    \midrule
    Baseline & -- & -- \\
    Intervention & -- & -- \\
    \bottomrule
  \end{tabular}
\end{table}

\section{Error Analysis and Discussion}
What fails, for whom, and under what conditions? Distinguish observations
from explanations and causal claims.

\section{Limitations and Responsible Use}
Describe what the experiment cannot establish, data and compute limits,
potential harms, and relevant licensing or privacy constraints.

\section{Conclusion}
Summarize only what the evidence supports and name the next falsifiable test.

\bibliographystyle{plainnat}
\bibliography{references}
\end{document}
